\documentclass[a4paper,12pt]{amsart}

\usepackage[foot]{amsaddr}

\usepackage{hyperref}
\hypersetup{
    colorlinks=true,
    urlcolor=blue,
    linkcolor=blue,
}

\usepackage{geometry}
\geometry{a4paper,margin=1in}

\usepackage{pdflscape}
\usepackage{longtable}
\usepackage{array}
\usepackage{booktabs}
\usepackage{diagbox}
\setlength{\LTcapwidth}{\textwidth}

\usepackage{pgfgantt}

\begin{document}

% Top matter
\title[Project Plan]{Project Plan for the Course COMP1110 -- Group D08}
% Authors, in alphabetical order by surname
%%
\newcommand{\cBanh}{C. Banh}
\author[\cBanh]{Chun Hei Banh}
\address{\cBanh, [Programme], [UID], \href{mailto:??@connect.hku.hk}{\texttt{??@connect.hku.hk}}}
%
\newcommand{\sLiu}{S. Liu}
\author[\sLiu]{Sze Pok Liu}
\address{\sLiu, [Programme], [UID], \href{mailto:u3658945@connect.hku.hk}{\texttt{u3658945@connect.hku.hk}}}
%
\newcommand{\jShing}{J. Shing}
\author[\jShing]{Zhan Ho Jacob Shing}
\address{\jShing, BEng(CompSc), 3036228892, \href{mailto:jacobszh@connect.hku.hk}{\texttt{jacobszh@connect.hku.hk}}}
%
\newcommand{\yYe}{Y. Ye}
\author[\yYe]{Yubing Ye}
\address{\yYe, [Programme], [UID], \href{mailto:u3629150@connect.hku.hk}{\texttt{u3629150@connect.hku.hk}}}
%
\newcommand{\yZhang}{Y. Zhang}
\author[\yZhang]{Yikun Zhang}
\address{\yZhang, BEng(CompSc), 3036482333, \href{mailto:pzdmmsd@connect.hku.hk}{\texttt{pzdmmsd@connect.hku.hk}}}
%%
\date{21 March 2026}
\maketitle

% BODY

\section{Topic Selection \& Introduction}

The group has selected the topic \textbf{B. Smart Public Transport Advisor} for the project.
For the project, the group will implement a route planner with the focus for pedestrian navigation within the campus of The University of Hong Kong (HKU).
To achieve that, the group will model the campus as an undirected multigraph as defined mathematically by:
\begin{equation} \label{eq:graph}
    G=(V,E),
\end{equation}
where for each vertex $v \in V$, $v$ represents a key location on the campus, such as a building entrance, a canteen, etc.;
and for each edge $e \in E$ represents an undirected path between two distinct vertices.
In addition to the graph structure, each edge will have a set of attributes, such as the average walking time, type of path (e.g., slope, stairs, etc.), etc.
The details of what attributes to include will be determined during the design phase of the project.

\section{Review of Existing Solutions}

For researching purposes, the group will review three types of existing solutions, each with two examples:
\begin{enumerate}
    \item \textbf{General purpose route planners}: Google Maps, Citymapper;
    \item \textbf{Transportation-specific route planners}: MTR Mobile App, HKBUS.app; and
    \item \textbf{Static navigation systems}: HKU campus map, On-campus navigation signages.
\end{enumerate}
The examples will be evaluated across several dimensions, such as easiness of use, specificity for pedestrian navigation, etc.
The conclusions drawn from the comparison shall be used to fine-tune the design of our own route planning implementation.

\section{Task Decomposition, Assignments \& Timeline}

The project will be decomposed into three phases: \textbf{Research \& Design}, \textbf{Implementation}, and \textbf{Test \& Evaluation}, each with its dedicated tasks.
Each task has a defined timeline and expected deliverables, which will be detailed in the following subsections.

\subsection{Task decomposition}
The tasks, along with their descriptions, expected deliverables, and time estimates, are listed in Table~\ref{tab:tasks}.

\begin{landscape}
\begin{longtable}{ll m{8cm} m{8cm} m{1cm}}
    \caption{A list of tasks for accomplishing the project, along with their descriptions, expected deliverables, and time estimations.}
    \label{tab:tasks}\\
    \toprule
    \textbf{Task ID} & \textbf{Phase} & \textbf{Description} & \textbf{Expected Deliverables} & \textbf{Time (days)} \\
    \midrule \endhead

    \texttt{R01}\label{task:r01} & Research \& Design & Survey existing navigation solutions and compare them in a table format. & Comparison table of existing tools (target users, key features, missing features, etc.) with summary takeaways. & 3 \\
    \midrule[0.2pt]
    \texttt{R02}\label{task:r02} & Research \& Design & Problem modeling for our solution: define core entities, input/output, and coverage scope. & Documentation of core concepts (stop, segment, journey, etc.), input/output specifications, scope of coverage. & 1 \\
    \midrule[0.2pt]
    \texttt{I01}\label{task:i01} & Implementation & Data collection: assign responsible areas to each team member, measure and record paths between designated nodes. & Collected path data (distance, time, cost, etc.), raw measurement records or data sheets. & 7 \\

    \bottomrule \endlastfoot
\end{longtable}
\end{landscape}

\subsection{Task assignments}
The task assignments is listed in Table~\ref{tab:task-assignments}.

\begin{table}
    \caption{The tasks as defined in Table~\ref{tab:tasks} and their assignees.}\label{tab:task-assignments}
    \begin{tabular}{ll}
        \toprule
        \textbf{Task ID} & \textbf{Assignee(s)} \\
        \midrule
        \hyperref[task:r01]{\texttt{R01}} & \cBanh, \sLiu, \jShing, \yYe, and \yZhang \\
        \hyperref[task:r02]{\texttt{R02}} & \dots \\
        \bottomrule
    \end{tabular}
\end{table}

\subsection{Plan timeline}
Note that the tasks set out in Table~\ref{tab:tasks} are not necessarily executed sequentially and linearly.
They may be arranged in parallel where possible and appropriate.
The timeline for the project and the dependency of tasks is illustrated by the Gantt chart in Figure~\ref{fig:project-gantt-chart}.

\begin{figure}
    Placeholder.
    \caption{The Gantt chart illustrating the project execution.}\label{fig:project-gantt-chart}
\end{figure}

\end{document}
